\chapter{Introduction}
\label{ch:introduction}

\section{Motivation}
\label{sec:intro-motivation}

The Guadalajara Metropolitan Area (\zmg) faces the mobility predicament common to
large Latin American cities: a spatially-expanding, car-dependent periphery served by a
fragmented, largely concessioned bus system, onto which a small premium tier of bus
rapid transit and light rail has recently been grafted. In this setting the central
planning question is not how to describe the network that exists, but where new transit
would do the most good---which corridors carry latent, unmet demand that the current
supply fails to serve. Getting this question wrong is expensive: transit investment is
large, spatially fixed, and effectively irreversible once built.

\section{The circularity problem in data-driven transit planning}
\label{sec:intro-circularity}

A growing body of work applies machine learning to transit-network questions, typically
by training a model to predict where stations or stops ``should'' be. A recurring
methodological flaw undermines much of it: the dependent variable is derived, directly
or indirectly, from where transit \emph{already exists}. A model trained to predict a
``has-a-stop'' label, or a suitability class that is itself a function of current
supply, can only learn to reproduce the status quo. Such a model may report near-perfect
accuracy while being, by construction, incapable of identifying unmet need---it rewards
the network that already exists rather than diagnosing where it is absent. Breaking this
circularity is the methodological starting point of this thesis.

\section{Research questions}
\label{sec:intro-rqs}

The thesis is organized around four research questions:

\begin{description}
  \item[RQ1 -- Demand.] Can transit demand be estimated at \ageb resolution from
    Tier-1 (non-transit) data alone, so that demand is never a function of current
    supply?
  \item[RQ2 -- Diagnosis.] Does an independent coverage-gap index, built from demand
    against supply, identify genuinely under-served areas---and can it be validated
    out-of-sample?
  \item[RQ3 -- Generation.] Can candidate corridors be generated and evaluated on
    genuine need, non-redundancy, and demand efficiency, rather than on agreement with
    the existing network?
  \item[RQ4 -- Transferability.] Does the framework transfer to other metropolitan
    areas without bespoke re-engineering?
\end{description}

\section{Contributions}
\label{sec:intro-contributions}

The thesis makes four contributions, matched to the research questions. First, a
Tier-1 transit-demand surface (\cref{sec:w1}) that estimates latent demand without any
transit input, calibrated against a metropolitan travel survey. Second, and most
significantly, a \emph{coverage-gap diagnostic} (\cref{sec:w3,sec:w4}) that replaces the
circular ``has-a-stop'' target with an independent demand-versus-supply gap, and that is
corroborated out-of-sample by the subsequent opening of \zmg Line~4. Third, a corridor
generator and multi-objective evaluation (\cref{sec:w5,sec:w6}) that scores candidate
lines on need, non-redundancy, and demand efficiency, and that---once corridors are
shaped as real paths and judged by anchor-directness---produces at least one
substantive, feasible, merit-passing corridor. Fourth, a demonstration that the entire
framework \emph{transfers} end-to-end to two further metropolitan areas of contrasting
size (\cref{ch:transferability}). A cross-cutting contribution is methodological
honesty: the thesis distinguishes what it validates strongly (the diagnostic) from what
it validates only partially (the generator), and argues why the two carry different
evidential weight (\cref{ch:discussion}).

\section{Thesis structure}
\label{sec:intro-structure}

The remainder of the thesis follows the pipeline. \Cref{ch:literature} reviews the
Node--Place tradition, demand modeling, multi-criteria weighting, and machine-learning
siting, and locates the research gap. \Cref{ch:data} describes the study area, the
\ageb unit, and the data. \Cref{ch:methodology} presents the eight-workstream
framework in full. \Cref{ch:results} reports its outputs for \zmg, and
\cref{ch:transferability} repeats the analysis for Toluca and Aguascalientes.
\Cref{ch:discussion} interprets the results and their limits, and
\cref{ch:conclusions} answers the research questions and outlines future work.
